% Conceptual Ashby material-property chart: Young's modulus E vs density rho.
% Material families as ellipse "bubbles" on log-log axes, with a constant
% stiffness-per-mass index guide line M1 = E^{1/2}/rho.
\documentclass{article}
\usepackage[paperwidth=120cm,paperheight=120cm,margin=10mm]{geometry}
\pagestyle{empty}
\usepackage{fontspec}
\usepackage[bidi=basic]{babel}
\babelprovide[import]{english}
\babelprovide[main, import]{persian}
\babelfont{rm}[Renderer=Node, Path=../fonts/, Extension=.ttf, UprightFont=*-Regular, BoldFont=*-Bold]{Vazirmatn}
\babelfont{sf}[Renderer=Node, Path=../fonts/, Extension=.ttf, UprightFont=*-Regular, BoldFont=*-Bold]{Vazirmatn}
\providecommand{\setRTL}{}\providecommand{\setLTR}{}
\providecommand{\RL}[1]{\foreignlanguage{persian}{#1}}
\providecommand{\LR}[1]{\babelsublr{#1}}
\usepackage{tikz}
\usepackage{pgfplots}
\pgfplotsset{compat=1.18}
\usetikzlibrary{arrows.meta, positioning, shapes.geometric, shapes.misc, calc, fit, backgrounds, decorations.pathmorphing, patterns, angles, quotes}
\begin{document}
\setRTL
\babelsublr{\begin{tikzpicture}[>=Stealth, font=\small]
  \begin{loglogaxis}[
      width=13cm, height=10cm,
      xlabel={چگالی $\rho_m$ \LR{(g/cm$^3$)}},
      ylabel={مدول یانگ $E$ (GPa)},
      xmin=1, xmax=20, ymin=1, ymax=600,
      xtick={1,2,4,8,16}, xticklabels={1,2,4,8,16},
      ytick={1,3,10,20,70,200,500},
      yticklabels={1,3,10,20,70,200,500},
      log ticks with fixed point,
      clip=false,
      every axis x label/.style={at={(ticklabel* cs:1.0)}, anchor=west},
      every axis y label/.style={at={(ticklabel* cs:1.0)}, anchor=south},
    ]

    % ---- material-family bubbles (ellipses in log space) ----
    % Ceramics (Al2O3, ZrO2): high E, moderate density
    \draw[fill=orange!20, draw=orange!70!black, thick] (axis cs:4.0,380) ellipse [x radius=8pt, y radius=20pt];
    \node[orange!60!black, align=center] at (axis cs:4.0,500) {سرامیک‌ها\\\scriptsize Al$_2$O$_3$, ZrO$_2$};

    % Metals region (steel 316L, Cu, Ag, Ti, Al)
    \draw[fill=blue!12, draw=blue!60!black, thick] (axis cs:6.5,120) ellipse [x radius=26pt, y radius=22pt];
    \node[blue!55!black] at (axis cs:8.0,193) {SS-316L};
    \node[blue!55!black] at (axis cs:4.43,114) {Ti};
    \node[blue!55!black] at (axis cs:2.70,69)  {Al};
    \node[blue!55!black] at (axis cs:9.0,100)  {Cu};

    % CFRP (composite, light, stiff)
    \draw[fill=green!18, draw=green!55!black, thick] (axis cs:1.6,130) ellipse [x radius=7pt, y radius=20pt];
    \node[green!45!black] at (axis cs:1.6,250) {CFRP};

    % GFRP / G10 composites
    \draw[fill=green!12, draw=green!45!black, thick] (axis cs:1.95,30) ellipse [x radius=8pt, y radius=10pt];
    \node[green!40!black] at (axis cs:1.95,12) {GFRP / G10};

    % Polymers (PEEK)
    \draw[fill=red!12, draw=red!55!black, thick] (axis cs:1.3,4) ellipse [x radius=7pt, y radius=10pt];
    \node[red!55!black] at (axis cs:1.3,1.7) {پلیمرها (PEEK)};

    % ---- material index guide line: E = C * rho^2  =>  E^{1/2}/rho = const ----
    % choose C so the line passes through (rho=2, E=20): C = 5
    \addplot[domain=1:20, samples=2, very thick, black, densely dashed] {5*x^2};
    \node[rotate=63, anchor=south, black] at (axis cs:6.2,250)
        {\footnotesize $M_1=E^{1/2}/\rho_m=\mathrm{const}$};

    % ---- improvement-direction arrow (up-left is better for M1) ----
    \draw[->, very thick, black] (axis cs:9,8) -- (axis cs:4.5,28)
        node[midway, sloped, above] {\footnotesize بهتر};

  \end{loglogaxis}
\end{tikzpicture}}
\end{document}
